% Part: model-theory
% Chapter: interpolation
% Section: interpolation-proof

\documentclass[../../../include/open-logic-section]{subfiles}

\begin{document}

\olfileid{mod}{int}{prf}

\olsection{Craig 内插定理}

\begin{thm}[Craig 内插定理]
\ollabel{thm:interpol}
如果 $\Entails !A \lif !B$，则存在一个语句 $!C$ 使得
$\Entails !A \lif !C$ 且 $\Entails !C \lif !B$，并且
$!C$ 中出现的每个常项符号、函数符号和谓词符号（$\eq$ 除外）
都同时出现在 $!A$ 与 $!B$ 中。语句 $!C$ 称为 $!A$ 与 $!B$ 的一个
\emph{内插项}。
\end{thm}

证明思路使用了``模型交错''技巧：给定 $\Entails !A \lif !B$
的一个反例模型，我们可以将其拆解并重新组合，以分别构造出
$\Entails !A \lif !C$ 和 $\Entails !C \lif !B$ 的反例，除非满足特定条件，
而该条件正好能给出所求的内插项 $!C$。

Craig 内插定理之所以有趣，部分原因在于它允许我们在理论之间进行``转译''。
例如，若语言 $L_1$ 中的 $!A$ 语义地后承语言 $L_2$ 中的 $!B$，
那么我们可以找到一个作为``桥梁''的 $!C$，它仅使用 $L_1$ 和 $L_2$ 共有的符号，
并且 $!A$ 语义地后承 $!C$，$!C$ 又语义地后承 $!B$。这样，
关于 $!A$ 和 $!B$ 的后承关系，就可以通过这个共同语言中的 $!C$ 来``解释''。

\begin{defn}[内插项]
\ollabel{defn:interpolant}
若语句 $!C$ 满足 \olref{thm:interpol} 的结论，则称 $!C$ 是
$!A$ 与 $!B$ 的一个\emph{内插项}。
\end{defn}

该定理的证明将在下一节中给出。

\begin{proof}
设 $\Lang{L}_1$ 是 $!A$ 的语言，$\Lang{L}_2$ 是 $!B$ 的语言。令 $\Lang{L}_0 = \Lang{L}_1 \cap \Lang{L}_2$。对每个 $i \in \{0, 1, 2 \}$，令 $\Lang{L}'_i$ 是通过向 $\Lang{L}_i$ 添加无穷多新的常项符号 $\Obj c_0,
\Obj c_1, \Obj c_2, \dots$ 而得到的语言。

如果 $!A$ 不可满足，则 $\lexists[x][\eq/[x][x]]$ 是一个内插项。如果 $\lnot !B$ 不可满足（因而 $!B$ 有效），则 $\lexists[x][\eq[x][x]]$ 是一个内插项。因此我们可假设 $!A$ 和 $\lnot !B$ 都是可满足的。

为证明内插定理的逆否命题，假设不存在 $!A$ 与 $!B$ 的内插项。换言之，假设 $\{!A\}$ 与 $\{\lnot !B\}$ 在 $\Lang{L}_0$ 中不可分。

我们的目标是将序对 $(\{ !A \}, \{\lnot!B\})$ 扩张为极大不可分序对 $(\Gamma^*, \Delta^*)$。令 $!A_0$, $!A_1$, $!A_2$, \dots 枚举 $\Lang{L}_1$ 的语句，$!B_0$, $!B_1$, $!B_2$, \dots 枚举 $\Lang{L}_2$ 的语句。我们如下定义两个递增的语句集序列 $(\Gamma_n, \Delta_n)$（$n \ge 0$）。令 $\Gamma_0 = \{ !A\}$ 且 $\Delta_0 = \{\lnot !B \}$。假设 $(\Gamma_n, \Delta_n)$ 已定义，按下述方式定义 $\Gamma_{n+1}$ 与 $\Delta_{n+1}$：
\begin{enumerate}
\item 如果 $\Gamma_n \cup \{!A_n \}$ 与 $\Delta_n$ 在 $\Lang{L}'_0$ 中不可分，则将 $!A_n$ 加入 $\Gamma_{n+1}$。此外，若 $!A_n$ 是存在公式 $\lexists[x][!S]$，则选取一个不在 $\Gamma_n$、$\Delta_n$、$!A_n$ 或 $!B_n$ 中出现的新常项符号 $c$，并将 $\Subst{!S}{c}{x}$ 加入 $\Gamma_{n+1}$。
\item 如果 $\Gamma_{n+1}$ 与 $\Delta_n \cup \{!B_n \}$ 在 $\Lang{L}'_0$ 中不可分，则将 $!B_n$ 加入 $\Delta_{n+1}$。此外，若 $!B_n$ 是存在公式 $\lexists[x][!S]$，则选取一个不在 $\Gamma_{n+1}$、$\Delta_n$、$!A_n$ 或 $!B_n$ 中出现的新常项符号 $c$，并将 $\Subst{!S}{c}{x}$ 加入 $\Delta_{n+1}$。
\end{enumerate}
最后，定义：
\begin{align*}
  \Gamma^* & = \bigcup_{n\ge 0} \Gamma_n, & 
  \Delta^* & = \bigcup_{n\ge 0} \Delta_n.
\end{align*}
现在我们可以通过对 $n$ 进行同时归纳来证明：
\begin{enumerate}
\item\ollabel{part-a} $\Gamma_n$ 与 $\Delta_n$ 在 $\Lang{L}'_0$ 中不可分；
\item\ollabel{part-b} $\Gamma_{n+1}$ 与 $\Delta_n$ 在 $\Lang{L}'_0$ 中不可分。
\end{enumerate}
\olref{part-a} 的归纳基础由 \olref[sep]{lem:sep1} 给出。对于 \olref{part-b}，我们需要区分三种情况：
\begin{enumerate}
\item 如果 $\Gamma_0 \cup \{!A_0 \}$ 与 $\Delta_0$ 可分，则 $\Gamma_1 = \Gamma_0$ 且 \olref{part-b} 就是 \olref{part-a}；
\item 如果 $\Gamma_1 = \Gamma_0 \cup\{ !A_0\}$，则根据构造，$\Gamma_1$ 与 $\Delta_0$ 不可分。
\item 剩下考虑 $!A_0$ 是存在公式的情形，即 $\Gamma_1 = \Gamma_0 \cup \{ \lexists[x][!S], \Subst{!S}{c}{x}
  \}$。根据构造，$\Gamma_0 \cup \{ \lexists[x][!S]\}$ 与 $\Delta_0$ 不可分，因此由 \olref[sep]{lem:sep2}，$\Gamma_0 \cup \{ \lexists[x][!S], \Subst{!S}{c}{x} \}$ 与 $\Delta_0$ 也不可分。
\end{enumerate}
这就完成了上述 \olref{part-a} 和 \olref{part-b} 的归纳基础。接下来考虑归纳步骤。对于 \olref{part-a}，如果 $\Delta_{n+1} = \Delta_n \cup \{ !B_n \}$，则根据构造（即使 $!B_n$ 是存在公式，由 \olref[sep]{lem:sep2} 亦保证），$\Gamma_{n+1}$ 与 $\Delta_{n+1}$ 不可分；如果 $\Delta_{n+1} = \Delta_n$（因为 $\Gamma_{n+1}$ 与 $\Delta_n \cup \{!B_n\}$ 可分），则我们运用对 \olref{part-b} 的归纳假设。对于 \olref{part-b} 的归纳步骤，如果 $\Gamma_{n+2} = \Gamma_{n+1} \cup
\{!A_{n+1} \}$，则根据构造（即使 $!A_{n+1}$ 是存在公式，由 \olref[sep]{lem:sep2} 亦保证），$\Gamma_{n+2}$ 与 $\Delta_{n+1}$ 不可分；如果 $\Gamma_{n+2} = \Gamma_{n+1}$，则我们运用刚刚证明的关于 \olref{part-a} 的归纳情形。这就完成了关于 \olref{part-a} 和 \olref{part-b} 的归纳。

由此可得 $\Gamma^*$ 与 $\Delta^*$ 是不可分的；若非如此，由紧致性，存在 $n \ge 0$ 使得 $\Gamma_n$ 与 $\Delta_n$ 可分，这与 \olref{part-a} 矛盾。特别地，$\Gamma^*$ 和 $\Delta^*$ 是一致的：因为如果前者或后者不一致，则它们将分别被 $\lexists[x][\eq/[x][x]]$ 或 $\lforall[x][\eq[x][x]]$ 所分离。

我们现在证明 $\Gamma^*$ 在 $\Lang{L}'_1$ 中是极大一致的，且类似地 $\Delta^*$ 在 $\Lang{L}'_2$ 中是极大一致的。对前者，假设对某个 $n \ge 0$，$!A_n \notin \Gamma^*$ 且 $\lnot !A_n \notin \Gamma^*$。若 $!A_n \notin \Gamma^*$，则 $\Gamma_n \cup \{!A_n \}$ 与 $\Delta_n$ 可分，于是存在 $!C \in \Lang{L}'_0$ 使得：
\begin{align*}
  \Gamma^* & \Entails !A_n \lif !C, & 
  \Delta^* & \Entails \lnot !C.
\end{align*}
类似地，若 $\lnot !A_n \notin \Gamma^*$，则存在 $!C' \in
\Lang{L}'_0$ 使得：
\begin{align*}
  \Gamma^* & \Entails \lnot !A_n \lif !C', & 
  \Delta^* & \Entails \lnot !C'.
\end{align*}
根据命题逻辑，$\Gamma^* \Entails !C \lor !C'$ 且 $\Delta^* \Entails \lnot (!C \lor !C')$，因此 $!C \lor !C'$ 分离了 $\Gamma^*$ 与 $\Delta^*$。类似的论证可确立 $\Delta^*$ 的极大性。

最后，我们证明 $\Gamma^* \cap \Delta^*$ 在 $\Lang{L}'_0$ 中是极大一致的。它显然是一致的，因为它是一致集合的交集。为证极大性，设 $!S \in \Lang{L}'_0$。现在，$\Gamma^*$ 在 $\Lang{L'_1}
\supseteq \Lang{L'_0}$ 中极大，类似地 $\Delta^*$ 在 $\Lang{L'_2} \supseteq \Lang{L'_0}$ 中极大。于是，要么 $!S \in \Gamma^*$ 要么 $\lnot !S \in \Gamma^*$，并且要么 $!S \in \Delta^*$ 要么 $\lnot !S \in \Delta^*$。若 $!S \in \Gamma^*$ 且 $\lnot !S \in \Delta^*$，则 $!S$ 将分离 $\Gamma^*$ 与 $\Delta^*$；若 $\lnot !S \in \Gamma^*$ 且 $!S \in \Delta^*$，则 $\Gamma^*$ 与 $\Delta^*$ 将被 $\lnot !S$ 分离。因此，要么 $!S \in \Gamma^* \cap \Delta^*$，要么 $\lnot !S \in \Gamma^* \cap \Delta^*$，故 $\Gamma^* \cap \Delta^*$ 是极大的。

由于 $\Gamma^*$ 极大一致，它有一个模型 $\Struct{M}'_1$，其论域 $\Domain{M'_1}$ 恰好由解释常项符号的元素 $\Assign{c}{M'_1}$ 组成——正如完备性定理（\olref[fol][com][cth]{thm:completeness}）的证明中那样。类似地，$\Delta^*$ 有一个模型 $\Struct{M}'_2$，其论域 $\Domain{M'_2}$ 由常项符号的解释 $\Assign{c}{M'_2}$ 给出。

令 $\Struct{M_1}$ 是通过从 $\Struct{M'_1}$ 中丢弃 $\Lang{L'_1} \setminus \Lang{L'_0}$ 中的常项符号、函数符号和谓词符号的解释而得到的，$\Struct{M_2}$ 类似地得到。那么由 $h(\Assign{c}{M'_1}) = \Assign{c}{M'_2}$ 定义的映射 $h \colon M_1 \to M_2$ 是 $\Lang{L}'_0$ 上的一个同构，因为如前所示，$\Gamma^* \cap \Delta^*$ 在 $\Lang{L}'_0$ 中极大一致。这之所以成立，是因为任何 $\Lang{L}'_0$-语句要么同时属于 $\Gamma^*$ 和 $\Delta^*$，要么都不属于：所以 $\Assign{c}{M'_1} \in
\Assign{P}{M'_1}$ 当且仅当 $\Atom{P}{c} \in \Gamma^*$ 当且仅当 $\Atom{P}{c} \in \Delta^*$ 当且仅当 $\Assign{c}{M'_2} \in
\Assign{P}{M'_2}$。同构满足的其他条件可以类似地确立。

现在为语言 $\Lang{L_1} \cup \Lang{L_2}$ 定义模型 $\Struct{M}$ 如下：
\begin{enumerate}
\item 论域 $\Domain{M}$ 就是 $\Domain{M_2}$，即所有 $\Assign{c}{M'_2}$ 的集合；
\item 若谓词符号 $P$ 在 $\Lang{L_2} \setminus
  \Lang{L_1}$ 中，则 $\Assign{P}{M} = \Assign{P}{M'_2}$；
\item 若谓词 $P$ 在 $\Lang{L}_1\setminus \Lang{L}_2$ 中，则 $\Assign{P}{M} = h(\Assign{P}{M'_2})$，即 $\tuple{\Assign{c_1}{M'_2}, \dots, \Assign{c_n}{M'_2}} \in
  \Assign{P}{M}$ 当且仅当 $\tuple{\Assign{c_1}{M'_1}, \dots,
  \Assign{c_n}{M'_1}} \in \Assign{P}{M'_1}$。
\item 若谓词符号 $P$ 在 $\Lang{L}_0$ 中，则 $\Assign{P}{M} =
  \Assign{P}{M'_2} = h(\Assign{P}{M'_1})$。
\item $\Lang{L}_1 \cup \Lang{L}_2$ 中的函数符号（包括常项符号）类似处理。
\end{enumerate}

最后，通过对公式归纳可以证明，$\Struct{M}$ 在所有 $\Lang{L'_1}$ 的公式上与 $\Struct{M'_1}$ 相符，在所有 $\Lang{L'_2}$ 的公式上与 $\Struct{M'_2}$ 相符。特别地，$\Struct{M} \Entails \Gamma^* \cup \Delta^*$，因此 $\Struct{M} \Entails !A$ 且 $\Struct{M} \Entails \lnot!B$，从而 $\not\Entails !A \lif !B$。这就完成了 Craig 内插定理的证明。
\end{proof}

\end{document}